\section{Reproducción de resultados}\label{app:comandos}

Todos los comandos se ejecutan desde la raíz del repositorio.

\begin{center}
\footnotesize
\begin{tabular}{ll}
  \toprule
  Resultado & Comando \\
  \midrule
  Coeficientes y secuencia QPSK & \texttt{python3 model/qpsk\_rrc.py --seed 1} \\
  Correspondencia entre ramas en flotante & \texttt{python3 -m tests.test\_a03\_time\_vs\_frequency} \\
  Barrido de SQNR y formatos & \texttt{python3 -m model.fixed\_point} \\
  Flotante contra fijo por caso & \texttt{python3 -m model.float\_vs\_fixed} \\
  Vectores de prueba & \texttt{python3 -m model.vectors} \\
  Integración: modelos y simulación & \texttt{make smoke} \\
  Regresión del banco de pruebas común & \texttt{make regresion} \\
  Integración completa con síntesis & \texttt{make verify} \\
  PPA SKY130, frecuencia & \texttt{sh scripts/ppa/openlane/run.sh --target frequency\_serial} \\
  PPA SKY130, tiempo & \texttt{sh scripts/ppa/openlane/run\_time.sh --target time\_serial} \\
  Tabla PPA & \texttt{python3 -m scripts.ppa\_table} \\
  \bottomrule
\end{tabular}
\end{center}

{\sloppy Las corridas SKY130 requieren OpenLane~2 instalado con Nix (\texttt{docs/frecuencia/ppa.md} y \texttt{docs/ppa-tiempo.md}); para las variantes paralelas se cambia el objetivo por \texttt{frequency\_optimized} o \texttt{time\_optimized}. La tabla PPA se genera a partir de la evidencia versionada en \texttt{reports/ppa/}, sin volver a correr las herramientas.\par}
